% GENERATED FILE. Edit canonical lessons, then run npm run generate:books.
% canonical-source-hash: fnv1a64:6537ed7c0057aaa2
% canonical-lessons: GE-C425-campingplatz, GE-C425-wohnwagen, GE-C425-wohnmobil, GE-C425-aufbauen, GE-C425-abbauen

\chapter{Auf dem Campingplatz}
\label{ch:ge-auf-dem-campingplatz}
\hlchaptermodality{425}

\begin{chapteropening}
\textbf{By the end of this chapter:} \emph{I can talk about camping: the campsite, the caravan, the camper van, setting things up and taking them down.}

Complete the last lesson of chapter 425: I can talk about camping: the campsite, the caravan, the camper van, setting things up and taking them down.
\end{chapteropening}

\section[der Campingplatz]{der Campingplatz — the campsite}

\label{lesson:GE-C425-campingplatz}



\begin{quote}
Before the new one: say the German for the overnight stay, then the German for the half board.
\end{quote}

\subsection*{You'll want to know: der Campingplatz}
\textbf{der Campingplatz} — \textquotedblleft{}the campsite\textquotedblright{}.

\begin{quote}
Der Campingplatz liegt direkt am Fluss.
\end{quote}

\begin{quote}
\emph{The campsite is right by the river.}
\end{quote}

Plural: \emph{die Campingplätze}.

\begin{cousinweb}
\emph{Campingplatz} is built from \emph{Camping}, borrowed from English \emph{camping}, and \emph{Platz} 'place', which came through Old French \emph{place} from Latin \emph{platea} (EtymDB). English \emph{camp} goes back to Latin \emph{campus} 'field' (Lewis \& Short).
\end{cousinweb}

\begin{grammarlens}[title={What you've built}]
One more word for telling, explaining and describing at B1.
\end{grammarlens}

\subsection*{Your turn}
\begin{itemize}
\raggedright
  \item \emph{Say it:} \emph{der Campingplatz}
  \item \emph{Say it:} \emph{der Campingplatz}, once more
  \item \emph{Say it:} a sentence of your own with \emph{der Campingplatz}
  \item \emph{Recall:} say the German for the overnight stay, then the German for the half board, then say \emph{der Campingplatz} again
\end{itemize}

\begin{tcolorbox}[breakable,colback=teal!4,colframe=teal!35!black,title={Before you move on}]
What does der Campingplatz mean? (\textquotedblleft{}The campsite\textquotedblright{}.) Say it once more.
\end{tcolorbox}
\section[der Wohnwagen]{der Wohnwagen — the caravan (US: trailer)}

\label{lesson:GE-C425-wohnwagen}



\begin{quote}
Before the new one: say the German for the half board, then the German for the campsite.
\end{quote}

\subsection*{You'll want to know: der Wohnwagen}
\textbf{der Wohnwagen} — \textquotedblleft{}the caravan (US: trailer)\textquotedblright{}.

\begin{quote}
Hinter dem Auto zieht er einen Wohnwagen.
\end{quote}

\begin{quote}
\emph{He is pulling a caravan behind the car.}
\end{quote}

Same form in the plural.

\begin{cousinweb}
\emph{Wohnwagen} is built from \emph{wohnen} 'to live', inherited from Old High German \emph{wonēn}, and \emph{Wagen} 'cart, vehicle', from Proto-Germanic \emph{wagnaz} (EtymDB).
\end{cousinweb}

\begin{grammarlens}[title={What you've built}]
One more word for telling, explaining and describing at B1.
\end{grammarlens}

\subsection*{Your turn}
\begin{itemize}
\raggedright
  \item \emph{Say it:} \emph{der Wohnwagen}
  \item \emph{Say it:} \emph{der Wohnwagen}, once more
  \item \emph{Say it:} a sentence of your own with \emph{der Wohnwagen}
  \item \emph{Recall:} say the German for the half board, then the German for the campsite, then say \emph{der Wohnwagen} again
\end{itemize}

\begin{tcolorbox}[breakable,colback=teal!4,colframe=teal!35!black,title={Before you move on}]
What does der Wohnwagen mean? (\textquotedblleft{}The caravan (US: trailer)\textquotedblright{}.) Say it once more.
\end{tcolorbox}
\section[das Wohnmobil]{das Wohnmobil — the camper van, the motorhome}

\label{lesson:GE-C425-wohnmobil}



\begin{quote}
Before the new one: say the German for the campsite, then the German for the caravan.
\end{quote}

\subsection*{You'll want to know: das Wohnmobil}
\textbf{das Wohnmobil} — \textquotedblleft{}the camper van, the motorhome\textquotedblright{}.

\begin{quote}
Mit dem Wohnmobil fahren sie drei Wochen durch das Land.
\end{quote}

\begin{quote}
\emph{They travel round the country for three weeks in the camper van.}
\end{quote}

Plural: \emph{die Wohnmobile}.

\begin{cousinweb}
\emph{Wohnmobil} is built from \emph{wohnen} 'to live' and \emph{mobil} 'mobile' (EtymDB), a learned borrowing from Latin \emph{mōbilis} 'movable' (Lewis \& Short).
\end{cousinweb}

\begin{grammarlens}[title={What you've built}]
One more word for telling, explaining and describing at B1.
\end{grammarlens}

\subsection*{Your turn}
\begin{itemize}
\raggedright
  \item \emph{Say it:} \emph{das Wohnmobil}
  \item \emph{Say it:} \emph{das Wohnmobil}, once more
  \item \emph{Say it:} a sentence of your own with \emph{das Wohnmobil}
  \item \emph{Recall:} say the German for the campsite, then the German for the caravan, then say \emph{das Wohnmobil} again
\end{itemize}

\begin{tcolorbox}[breakable,colback=teal!4,colframe=teal!35!black,title={Before you move on}]
What does das Wohnmobil mean? (\textquotedblleft{}The camper van, the motorhome\textquotedblright{}.) Say it once more.
\end{tcolorbox}
\section[aufbauen]{aufbauen — to put up, to set up}

\label{lesson:GE-C425-aufbauen}



\begin{quote}
Before the new one: say the German for the caravan, then the German for the camper van.
\end{quote}

\subsection*{You'll want to know: aufbauen}
\textbf{aufbauen} — \textquotedblleft{}to put up, to set up\textquotedblright{}.

\begin{quote}
Am Abend bauen sie einen Tisch und Stühle vor dem Wohnwagen auf.
\end{quote}

\begin{quote}
\emph{In the evening they set up a table and chairs in front of the caravan.}
\end{quote}

Separable: \emph{wir bauen auf}, \emph{wir haben aufgebaut}.

\begin{cousinweb}
\emph{aufbauen} is built from \emph{auf} 'up' and \emph{bauen} 'to build', inherited from Proto-Germanic \emph{būaną} (EtymDB).
\end{cousinweb}

\begin{grammarlens}[title={What you've built}]
One more word for telling, explaining and describing at B1.
\end{grammarlens}

\subsection*{Your turn}
\begin{itemize}
\raggedright
  \item \emph{Say it:} \emph{aufbauen}
  \item \emph{Say it:} \emph{aufbauen}, once more
  \item \emph{Say it:} a sentence of your own with \emph{aufbauen}
  \item \emph{Recall:} say the German for the caravan, then the German for the camper van, then say \emph{aufbauen} again
\end{itemize}

\begin{tcolorbox}[breakable,colback=teal!4,colframe=teal!35!black,title={Before you move on}]
What does aufbauen mean? (\textquotedblleft{}To put up, to set up\textquotedblright{}.) Say it once more.
\end{tcolorbox}
\section[abbauen]{abbauen — to take down, to dismantle; to reduce}

\label{lesson:GE-C425-abbauen}



\begin{quote}
Before the new one: say the German for the camper van, then the German for to put up.
\end{quote}

\subsection*{You'll want to know: abbauen}
\textbf{abbauen} — \textquotedblleft{}to take down, to dismantle; to reduce\textquotedblright{}.

\begin{quote}
Am letzten Tag bauen wir alles wieder ab.
\end{quote}

\begin{quote}
\emph{On the last day we take everything down again.}
\end{quote}

Separable: \emph{wir bauen ab}. \emph{Stress abbauen} — to reduce stress.

\begin{cousinweb}
\emph{abbauen} is built from \emph{ab} 'off, down' and \emph{bauen} 'to build', inherited from Proto-Germanic \emph{būaną} (EtymDB).
\end{cousinweb}

\begin{grammarlens}[title={What you've built}]
One more word for telling, explaining and describing at B1.
\end{grammarlens}

\subsection*{Your turn}
\begin{itemize}
\raggedright
  \item \emph{Say it:} \emph{abbauen}
  \item \emph{Say it:} \emph{abbauen}, once more
  \item \emph{Say it:} a sentence of your own with \emph{abbauen}
  \item \emph{Recall:} say the German for the camper van, then the German for to put up, then say \emph{abbauen} again
\end{itemize}

\begin{tcolorbox}[breakable,colback=teal!4,colframe=teal!35!black,title={Before you move on}]
What does abbauen mean? (\textquotedblleft{}To take down, to dismantle; to reduce\textquotedblright{}.) Say it once more.
\end{tcolorbox}
